%%%%%%%%%%%%%%%%%%%%%%%%%%%%%%%%%%%%%%%%%%%%%%%%%%%%%%%%%%%%%%%%%%%%%%%%%%%%%%%%
%% Capitulo 6: Plano e Especificação de Testes (STP / STD)
%% Sistema: PostoSmart - Automação e Gestão de Postos de Combustíveis
%% Autor: Diogo Santos Pires Jandiroba
%%%%%%%%%%%%%%%%%%%%%%%%%%%%%%%%%%%%%%%%%%%%%%%%%%%%%%%%%%%%%%%%%%%%%%%%%%%%%%%%

\chapter{Plano e Especificação de Testes de Software}
\label{cap:plano_testes}

\section{Estratégia e Pirâmide de Testes}

A estratégia de garantia de qualidade funcional e de integração do PostoSmart está estruturada na pirâmide de testes da \autoref{fig:piramide_testes_posto}, em estrita conformidade com a norma **ISO/IEC/IEEE 29119:2022**.

\begin{figure}[htbp]
\centering
\begin{tikzpicture}[every node/.style={font=\small\sffamily}]
    % Trapézios da Pirâmide
    \draw[thick, primary, fill=bgneutral] (-1.5, 3) -- (1.5, 3) -- (2.5, 1.8) -- (-2.5, 1.8) -- cycle;
    \node[align=center] at (0, 2.4) {\textbf{Testes E2E e UI}\\10\% (Cenários de Caixa e SEFAZ)};
    
    \draw[thick, primary, fill=borderline!40] (-2.5, 1.8) -- (2.5, 1.8) -- (3.8, 0.6) -- (-3.8, 0.6) -- cycle;
    \node[align=center] at (0, 1.2) {\textbf{Testes de Integração}\\30\% (Mock de Concentrador RS-485 / SQLite)};
    
    \draw[thick, primary, fill=white] (-3.8, 0.6) -- (3.8, 0.6) -- (5.2, -0.6) -- (-5.2, -0.6) -- cycle;
    \node[align=center] at (0, 0.0) {\textbf{Testes Unitários Automatizados}\\60\% (Regras de Bloqueio, Cálculos e Parsers)};
\end{tikzpicture}
\caption{Pirâmide de Testes de Software do PostoSmart (ISO 29119)}\label{fig:piramide_testes_posto}
\end{figure}

\section{Casos de Teste Padronizados}

A \autoref{tab:casos_teste_posto} especifica os cenários de verificação com entradas, ações e resultados esperados.

\begin{table}[htbp]
\caption{Especificação Detalhada de Casos de Teste (ISO/IEC/IEEE 29119)}
\label{tab:casos_teste_posto}
\centering
\small
\begin{tabularx}{\textwidth}{@{} l l Y Y c @{}}
\toprule
\textbf{ID} & \textbf{Requisito} & \textbf{Procedimento de Execução} & \textbf{Resultado Esperado} & \textbf{Resultado} \\
\midrule
CT01 & RF01, RN01 & Desengatar bico 01 sem pendências e aguardar pulso. & Concentrador autoriza fluxo; sistema marca bico como \texttt{ABASTECENDO}. & \badgeconcluido \\
CT02 & RF01, RN01 & Gerar 2 abastecimentos pendentes no bico 02 e tentar novo desengate. & Sistema bloqueia autorização (RN01) e aciona alarme no PDV. & \badgeconcluido \\
CT03 & RF03, RN03 & Simular queda da SEFAZ com mock de timeout ($> 3\,\text{s}$) no fechamento. & Sistema chaveia para contingência offline e imprime DANFE NFC-e. & \badgeconcluido \\
CT04 & RNF01 & Disparar telemetria simultânea em 16 bicos em loop a $50\,\text{ms}$. & Latência de atualização da tela do PDV inferior a $200\,\text{ms}$ (P99). & \badgeconcluido \\
\bottomrule
\end{tabularx}
\end{table}

\section{Matriz de Rastreabilidade (Requisitos \texorpdfstring{$\times$}{x} Testes)}

A \autoref{tab:rastreabilidade_testes_posto} confirma a cobertura total entre requisitos de negócio e baterias de teste.

\begin{table}[htbp]
\caption{Matriz de Rastreabilidade Requisitos $\times$ Casos de Teste}
\label{tab:rastreabilidade_testes_posto}
\centering
\small
\begin{tabularx}{\textwidth}{@{} l l l c @{}}
\toprule
\textbf{Requisito} & \textbf{Descrição Sintética} & \textbf{Casos de Teste Vinculados} & \textbf{Status} \\
\midrule
RF01 & Captura de Telemetria de Bico & CT01, CT02, CT04 & \badgeconcluido \\
RF02 & Fechamento de Venda no PDV & CT01, CT03 & \badgeconcluido \\
RF03 & Emissão de Documento Fiscal & CT03 & \badgeconcluido \\
RN01 & Bloqueio por Dupla Pendência & CT02 & \badgeconcluido \\
RNF01 & Latência de Telemetria ($\le 200\,\text{ms}$) & CT04 & \badgeconcluido \\
\bottomrule
\end{tabularx}
\end{table}
